%% kapa: exam=1 round=33 year=2022 period=2 form=A subject=appraisal_law subject_ko=감정평가관계법규 seq=14 number=14
%% source: Q-Net 감정평가사 기출문제 (cst00309) artlSeq=5209997 file=raw/1008/5209997/2022년도 제33회 감정평가사 1차시험 2교시(감정평가 관계법규, 회계학).pdf
%% conversion: kordoc 4.24.51 -> tools/gen1.py (auto); review=proofread-sonnet; flags=none
\begin{question}{14}
감정평가 및 감정평가사에 관한 법령상 감정평가법인등이 토지를 감정평가하는 경우 해당 토지의 임대료, 조성비용 등을 고려하여 감정평가를 할 수 있는 경우가 아닌 것은?
\begin{choices}
  \item 보험회사의 의뢰에 따른 감정평가
  \item 신탁회사의 의뢰에 따른 감정평가
  \item 「자산재평가법」에 따른 감정평가
  \item 법원에 계속 중인 소송을 위한 감정평가 중 보상과 관련된 감정평가
  \item 금융기관의 의뢰에 따른 감정평가
\end{choices}
\end{question}
\answer{4}
